\Needspace{7\baselineskip}\Article{303}
Die nach ihren verschiedenen Klassifikationen geordnete Tafel der negativen Determinanten bietet viele weitere bemerkenswerte Beobachtungen. Bei Determinanten der Form \(-(8n+3)\) ist die Anzahl der Klassen, sowohl in allen eigentlich primitiven Geschlechtern zusammen als auch in jedem einzelnen, stets durch \(3\) teilbar, mit der einzigen Ausnahme der Determinante \(-3\). Der Grund folgt unmittelbar aus Art.\ 256, VI. Bei Determinanten, deren Formen nur ein Geschlecht bilden, ist die Anzahl der Klassen stets ungerade. Für eine solche Determinante gibt es nämlich nur eine ambige Klasse, die Hauptklasse. Die übrigen Klassen treten paarweise entgegengesetzt auf; ihre Anzahl muss also gerade und damit die Gesamtzahl ungerade sein. Diese letzte Eigenschaft gilt auch für positive Determinanten. — Ferner scheint die Reihe der Determinanten, denen dieselbe gegebene Klassifikation entspricht, also eine gegebene Anzahl sowohl der Geschlechter als auch der Klassen, stets abzubrechen. Diese ziemlich überraschende Beobachtung erläutern wir an einigen Beispielen. Die erste, römische Zahl bezeichnet die Anzahl der eigentlich primitiven Geschlechter, bei negativer Determinante der positiven; die folgende die Anzahl der Klassen in jedem Geschlecht. Danach folgt die Reihe der Determinanten mit dieser Klassifikation; ihr negatives Vorzeichen ist der Kürze wegen weggelassen.

\SourceLeaf{372}

\begin{center}\begin{tabularx}{\textwidth}{@{}r@{\quad}X@{}}
I. 1\;\ldots & 1, 2, 3, 4, 7\\[2pt]
I. 3\;\ldots & 11, 19, 23, 27, 31, 43, 67, 163\\[2pt]
I. 5\;\ldots & 47, 79, 103. 127\\[2pt]
I. 7\;\ldots & 71, 151, 223, 343, 463, 487\\[2pt]
II. 1\;\ldots & 5, 6, 8, 9, 10, 12, 13, 15, 16, 18, 22, 25, 28, 37, 58\\[2pt]
II. 2\;\ldots & 14, 17, 20, 32, 34, 36, 39, 46, 49, 52, 55, 63, 64, 73, 82, 97, 100, 142, 148, 193\\[2pt]
IV. 1\;\ldots & 21, 24, 30, 33, 40, 42, 45, 48, 57, 60, 70, 72, 78, 85, 88, 93, 102, 112, 130, 133, 177, 190, 232, 253\\[2pt]
VIII. 1\;\ldots & 105, 120, 165, 168, 210, 240, 273, 280, 312, 330, 345. 357, 385, 408, 462, 520, 760\\[2pt]
XVI. 1\;\ldots & 840, 1320, 1365, 1848\\[2pt]
\end{tabularx}\end{center}
Ebenso findet man \(20\) Determinanten, die größte \(=-1423\), mit der Klassifikation I.~9; \(4\), die größte \(=-1303\), mit der Klassifikation I.~11 usw.\ Den Klassifikationen II.~3; II.~4; II.~5; IV.~2 entsprechen der Reihe nach nicht mehr als \(48,31,44,69\) Determinanten, deren größte \(-652,-862,-1318\), \(-1012\) sind. Die Tafel, der wir diese Beispiele entnommen haben, ist weit über die größten hier vorkommenden Determinanten hinaus fortgesetzt worden\HistoricalNote{\ensuremath{\ast})}{Während des Drucks ist sie \textit{ohne Unterbrechung} bis \(-3000\) fortgesetzt, außerdem durch den ganzen zehnten Tausender und mehrere andere verstreute Hunderter, dazu sehr viele sorgfältig ausgewählte einzelne Determinanten.}, ohne dass weitere zu diesen Klassifikationen gehörige Determinanten aufgetreten wären. Daher scheint kein Zweifel, dass die angegebenen Reihen wirklich abgebrochen sind, und nach Analogie darf man denselben Schluss auf jede andere Klassifikation ausdehnen. Da sich zum Beispiel im ganzen zehnten Tausender der Determinanten keine mit weniger als \(24\) Klassen gezeigt hat, ist es höchst wahrscheinlich, dass die Klassifikationen I.~23; I.~21 usw.; II.~11; II.~10 usw.; IV.~5; IV.~4; IV.~3; VIII.~2 schon vor \(-9000\) aufgehört haben oder wenigstens nur sehr wenigen Determinanten jenseits von \(-10000\) zukommen. \textit{Strenge} Beweise dieser Beobachtungen scheinen aber sehr schwierig. — Nicht weniger bemerkenswert ist, dass alle Determinanten, deren Formen sich auf \(32\) oder mehr Geschlechter verteilen, mindestens zwei Klassen in jedem Geschlecht haben. Die Klassifikationen XXXII.~1, LXIV. 1 usw.\ fallen also ganz weg; der kleinsten solchen Determinante, \(-9240\), entspricht XXXII.~2. Es erscheint ziemlich wahrscheinlich, dass mit wachsender Anzahl der Geschlechter immer mehr Klassifikationen wegfallen. In dieser Hinsicht sind die oben angegebenen \(65\) Determinanten,
\SourceLeaf{373}
denen die Klassifikationen I.~1; II.~1; IV.~1; VIII.~1; XVI.~1 entsprechen, sehr bemerkenswert. Man sieht leicht, dass sie alle und nur sie die beiden besonderen Eigenschaften haben, dass sämtliche zugehörigen Formenklassen ambig sind und beliebige Formen desselben Geschlechts notwendig sowohl eigentlich als auch uneigentlich äquivalent sind. Dieselben \(65\) Zahlen hat übrigens schon der berühmte Euler unter einem etwas anderen, unten zu erwähnenden Gesichtspunkt und mit einem leicht zu beweisenden Kriterium angegeben: \textit{Nouv. Mém. de l’Ac. de Berlin}, 1776, S.\ 338.

\Needspace{7\baselineskip}\Article{304}
Die Anzahl der eigentlich primitiven Klassen binärer Formen mit einer positiven Determinante, die eine \textit{Quadratzahl} \(kk\) ist, lässt sich vollständig im Voraus angeben. Sie ist gleich der Anzahl der zu \(2k\) teilerfremden und darunterliegenden Zahlen. Daraus folgt durch nicht schwierige, hier aber zu übergehende Überlegungen, dass sich die durchschnittliche Anzahl der Klassen für solche Determinanten um \(kk\) näherungsweise durch \(\frac{8k}{\pi\pi}\) ausdrückt. — Positive nichtquadratische Determinanten zeigen in dieser Hinsicht aber ganz besondere Erscheinungen. Eine kleine Klassenanzahl, zum Beispiel die Klassifikation I.~1 oder I.~3 oder II.~1 usw., tritt bei negativen und quadratischen Determinanten nur für kleine Werte auf und hört bald ganz auf. Dagegen liefert der weitaus größte Teil der positiven nichtquadratischen Determinanten, wenigstens der nicht sehr großen, Klassifikationen mit nur einer Klasse in jedem Geschlecht, sodass I.~3; I.~5; II.~2; II.~3; IV.~2 usw.\ sehr selten sind. So findet man zum Beispiel unter den \(90\) nichtquadratischen Determinanten unter \(100\) deren \(11,48,27\) mit den Klassifikationen I.~1, II.~1, IV.~1 der Reihe nach. Nur eine, \((37)\), hat I.~3; zwei, \((34\text{ und }82)\), haben II.~2; eine, \((79)\), hat II.~3. Mit wachsenden Determinanten werden größere Klassenanzahlen aber allmählich häufiger. Unter den \(96\) nichtquadratischen Determinanten von \(101\) bis \(200\) haben zwei, \((101,197)\), I.~3; vier, \((145,146,178,194)\), II.~2; drei, \((141,148,189)\), II.~3. Von den \(197\) Determinanten von \(801\) bis \(1000\) haben drei I.~3; vier II.~2; vierzehn II.~3; zwei II.~5; zwei II.~6; fünfzehn IV.~2; sechs IV.~3; zwei IV.~4; vier VIII.~2; die übrigen \(145\) haben eine Klasse in jedem Geschlecht. — Es wäre eine interessante, des Scharfsinns der Mathematiker würdige Frage, nach welchem Gesetz die Determinanten mit einer Klasse in jedem Geschlecht immer seltener werden. Bisher können wir weder theoretisch entscheiden noch aus Beobachtungen mit hinreichender Sicherheit vermuten, ob sie schließlich ganz abbrechen, was allerdings wenig wahrscheinlich erscheint, oder wenigstens \textit{unendlich selten} werden oder ob ihre Häufigkeit sich immer mehr
\SourceLeaf{374}
einer festen Grenze nähert. Die durchschnittliche Anzahl der Klassen wächst nur wenig schneller als die Anzahl der Geschlechter und weit langsamer als die Quadratwurzeln der Determinanten. Zwischen \(800\) und \(1000\) findet man sie \(=5,01\). Diesen Beobachtungen sei eine weitere hinzugefügt, die die Analogie zwischen positiven und negativen Determinanten gewissermaßen wiederherstellt. Wir haben nämlich gefunden, dass bei positiver Determinante \(D\) weniger die Klassenanzahl selbst als vielmehr ihr Produkt mit dem Logarithmus von \(t+u\surd D\), wobei \(t,u\) die kleinsten, von \(1,0\) verschiedenen Lösungen der Gleichung \(tt-Duu=1\) sind, der Klassenanzahl bei negativer Determinante in mehreren, hier nicht näher auszuführenden Beziehungen entspricht. Der Durchschnittswert dieses Produkts wird ebenso genau durch eine Formel der Gestalt \(m\surd D-n\) ausgedrückt. Die Werte der Konstanten \(m,n\) konnten wir bisher aber nicht theoretisch bestimmen. Soweit man aus dem Vergleich einiger Hunderter von Determinanten schließen darf, scheint \(m\) wenig von \(2\frac13\) abzuweichen. — Über die Grundsätze der vorigen Untersuchungen zu Durchschnittswerten von Größen, die keinem analytischen Gesetz folgen, sich einem solchen aber asymptotisch immer mehr nähern, behalten wir uns eine ausführlichere Behandlung bei anderer Gelegenheit vor. Wir gehen nun zu einer anderen Untersuchung über, die verschiedene eigentlich primitive Klassen derselben Determinante miteinander vergleicht und diesen langen Abschnitt abschließt.

\Needspace{10\baselineskip}\begin{center}\textit{Ein besonderer Algorithmus für eigentlich primitive Klassen: reguläre und irreguläre Determinanten usw.}\end{center}

\Needspace{7\baselineskip}\Article{305}
\textsc{Satz.} \textit{Es bezeichne \(K\) die Hauptklasse der Formen einer gegebenen Determinante \(D\) und \(C\) eine beliebige andere Klasse im Hauptgeschlecht der Formen derselben Determinante. Ferner seien \(2C,3C, 4C\) usw.\ die Klassen, die durch Verdopplung, Verdreifachung, Vervierfachung usw.\ der Klasse \(C\) entstehen, wie in Art.\ \(249\). Setzt man die Folge \(C,2C,3C\) usw.\ weit genug fort, so erreicht man schließlich eine mit \(K\) identische Klasse. Ist \(mC\) die erste mit \(K\) identische Klasse und die Gesamtzahl der Klassen im Hauptgeschlecht \(=n\), so gilt entweder \(m=n\), oder \(m\) ist ein echter Teiler von \(n\).}

\textit{Beweis.} I.\ Alle Klassen \(K,C,2C,3C\) usw.\ gehören notwendig zum Hauptgeschlecht, nach Art.\ \(247\). Daher können die ersten \(n+1\) Klassen dieser Reihe, \(K,C,2C\ldots nC\), offenbar nicht alle verschieden sein. Also ist entweder \(K\) mit einer der Klassen \(C,2C,3C\ldots nC\) identisch, oder wenigstens zwei dieser Klassen sind miteinander identisch.
\SourceLeaf{375}
Es sei \(rC=sC\) und \(r>s\). Dann gilt auch
\[\begin{gathered}(r-1)C=(s-1)C,\quad(r-2)C=(s-2)C\\\text{usw. und }(r+1-s)C=C\end{gathered}\]
also \((r-s)C=K\). \textit{Damit ist die erste Behauptung bewiesen.}

II.\ Daraus folgt auch unmittelbar entweder \(m=n\) oder \(m<n\). Es bleibt nur zu zeigen, dass im zweiten Fall \(m\) ein echter Teiler von \(n\) ist. Da die Klassen
\[K,C,2C\ldots(m-1)C,\text{ deren Gesamtheit wir mit }\mathfrak C\]
bezeichnen, in diesem Fall das Hauptgeschlecht noch nicht ausschöpfen, sei \(C'\) eine Klasse dieses Geschlechts, die nicht in \(\mathfrak C\) enthalten ist. Wir betrachten die Gesamtheit der Klassen, die durch Zusammensetzung von \(C'\) mit den einzelnen Klassen in \(\mathfrak C\) entstehen, nämlich
\[C',C'+C,C'+2C\ldots C'+(m-1)C\text{, mit der Bezeichnung }\mathfrak C'\]
Offenbar sind alle Klassen in \(\mathfrak C'\) sowohl voneinander als auch von allen in \(\mathfrak C\) verschieden und gehören zum Hauptgeschlecht. Schöpfen also \(\mathfrak C\) und \(\mathfrak C'\) dieses Geschlecht aus, so gilt \(n=2m\), andernfalls \(2m<n\). Im zweiten Fall sei \(C''\) eine Klasse des Hauptgeschlechts, die weder in \(\mathfrak C\) noch in \(\mathfrak C'\) enthalten ist. Man betrachte die Gesamtheit der Klassen, die durch Zusammensetzung von \(C''\) mit den einzelnen Klassen in \(\mathfrak C\) entstehen, also
\[C'',C''+C,C''+2C\ldots C''+(m-1)C\text{, mit der Bezeichnung }\mathfrak C''\]
Diese sind offenbar alle voneinander und von allen in \(\mathfrak C\) und \(\mathfrak C'\) verschieden und gehören zum Hauptgeschlecht. Wenn daher \(\mathfrak C,\mathfrak C',\mathfrak C''\) dieses Geschlecht ausschöpfen, ist \(n=3m\), andernfalls \(n>3m\). Dann zeigt eine ebenso behandelte weitere Klasse \(C'''\) des Hauptgeschlechts, die weder in \(\mathfrak C,\mathfrak C'\) noch in \(\mathfrak C''\) enthalten ist, dass entweder \(n=4m\) oder \(n>4m\) gilt, und so weiter. Da \(n\) und \(m\) endliche Zahlen sind, wird das Hauptgeschlecht schließlich notwendig ausgeschöpft. Dann ist \(n\) ein Vielfaches von \(m\), oder \(m\) ein echter Teiler von \(n\). \textit{Damit ist die zweite Behauptung bewiesen.}

\begingroup\rightskip=0pt plus 3em\relax
\textit{Beispiel.} Es sei \(D=-356,\ C=(5,2,72)\)\HistoricalNote{\ensuremath{\ast})}{Klassen werden hier stets durch die in ihnen enthaltenen einfachsten \textit{Formen} bezeichnet.}. Dann findet man \(2C=(20,8,21), 3C=(4,0,89),\ 4C=(20,-8,21),\ 5C=(5,-2,72),\ 6C=(1,0,356).\) Hier
\SourceLeaf{376}
ist also \(m=6\), während \(n\) bei dieser Determinante \(12\) ist. Wählt man für \(C'\) die Klasse \((8,2,45)\), so sind die übrigen fünf Klassen in \(\mathfrak C'\) gleich \((9,-2,40)\), \((9,2,40)\), \((8,-2,45)\), \((17,1,21)\), \((17,-1,21)\).
\par\endgroup
